\documentclass[11pt,letterpaper]{article}
\usepackage[margin=0.85in,headheight=16pt]{geometry}
\usepackage[T1]{fontenc}
\usepackage[utf8]{inputenc}
\usepackage{lmodern,microtype}
\usepackage{amsmath,amssymb,mathtools,booktabs,longtable,tabularx,array}
\usepackage[dvipsnames]{xcolor}
\usepackage{hyperref,fancyhdr,titlesec,enumitem,graphicx,tcolorbox}
\usepackage{lastpage}
\definecolor{ink}{HTML}{14263B}
\definecolor{sea}{HTML}{167C8A}
\definecolor{gold}{HTML}{C49747}
\definecolor{pale}{HTML}{EEF5F6}
\definecolor{graytext}{HTML}{516275}
\hypersetup{colorlinks=true,linkcolor=ink,urlcolor=sea,bookmarksopen=true,pdfauthor={Independent Educational Engineering Suite}}
\titleformat{\section}{\Large\bfseries\color{ink}}{\thesection}{0.7em}{}[\vspace{-0.3em}\color{gold}\titlerule]
\titleformat{\subsection}{\large\bfseries\color{sea}}{\thesubsection}{0.6em}{}
\setlist[itemize]{leftmargin=1.3em,itemsep=2pt,topsep=3pt}
\setlist[enumerate]{leftmargin=1.5em,itemsep=2pt,topsep=3pt}
\pagestyle{fancy}
\fancyhf{}
\fancyhead[L]{\small\sffamily\color{graytext}SECURITY ANALYSIS / ENGINEERING SUITE}
\fancyhead[R]{\small\sffamily\color{graytext}\textbf{\thepage}/\pageref{LastPage}}
\renewcommand{\headrulewidth}{0.3pt}
\renewcommand{\footrulewidth}{0pt}
\newtcolorbox{methodbox}[1][]{colback=pale,colframe=sea,boxrule=0.55pt,arc=1mm,title={\bfseries Method note},fonttitle=\sffamily,#1}
\newtcolorbox{sourcebox}[1][]{colback=white,colframe=gold,boxrule=0.55pt,arc=1mm,title={\bfseries Source boundary},fonttitle=\sffamily,#1}
\newcommand{\term}[1]{\textbf{\color{ink}#1}}
\newcommand{\src}{\textsf{\color{sea}[SOURCE]}}
\newcommand{\ext}{\textsf{\color{gold}[EXTENSION]}}
\newcommand{\gap}{\textsf{\color{red}[SOURCE GAP]}}
\newcommand{\answer}[1]{\par\smallskip\noindent\textbf{Check:} #1\par}
\newcommand{\cover}[3]{%
\begin{titlepage}\pagecolor{ink}\color{white}\vspace*{1.15in}
{\sffamily\large\color{gold}INDEPENDENT TECHNICAL DOCUMENTATION SUITE\par}\vspace{0.4in}
{\Huge\sffamily\bfseries #1\par}\vspace{0.15in}
{\Large\sffamily #2\par}\vspace{0.65in}
\color{gold}\rule{\textwidth}{2pt}\par\vspace{0.4in}
{\large #3\par}\vfill
{\small\sffamily Based on a critical reading of Graham and Dodd's \emph{Security Analysis}, Seventh Edition (2023). Original educational analysis; not affiliated with the publisher.\par}
\end{titlepage}\nopagecolor\color{black}}
\newcommand{\front}[3]{\cover{#1}{#2}{#3}\tableofcontents\clearpage
\begin{sourcebox}\src\ means a concept is attributable to the relevant chapter or seventh-edition contributor, summarized independently rather than reproduced. \ext\ means a modern method or engineering recommendation prepared for this suite, not an original Graham--Dodd rule. Historical examples and standards should not be treated as current market requirements. All numerical examples are synthetic, educational, and not securities recommendations.\end{sourcebox}}
\setlength{\parskip}{5pt}\setlength{\parindent}{0pt}
\begin{document}
\front{Portfolio Risk}{Engineering Policy, Metrics and Controls}{Permanent-loss focus with risk taxonomy, exposure limits, stress, liquidity and governance}

\section{Policy objectives and risk philosophy}
\src\ The historical emphasis on principal safety, adverse-condition analysis, permanent loss, enforceable claims, and continuous supervision provides the conceptual foundation. \ext\ Portfolio construction, risk statistics, limits, model governance and escalation procedures below are independent contemporary methods.

\term{Primary objective:} reduce the probability and severity of permanent impairment and forced liquidation while meeting documented mandate, liquidity and spending constraints. Volatility is relevant but not an adequate stand-alone proxy for loss of capital.

\section{Governance and the investment policy statement}
Document target return, allowed instruments, benchmark use, horizon, leverage, maximum drawdown tolerance, spending/withdrawal policy, daily cash buffer, concentration and issuer limits, currency policy, valuation rules, restricted assets, escalation powers, and review cadence. Risk must have authority to halt exposures breaching hard constraints; approvals must not depend solely on model-scored comfort.

\section{Risk taxonomy}
\begin{center}\small\begin{tabularx}{\textwidth}{p{0.2\textwidth}X X}\toprule Risk & Representative measurement & Control action \\
\midrule Permanent loss & Recovery downside, equity impairment & Claim/risk-capital gate \\
Market/valuation & Sensitivities, factor beta, drawdown & Scenario limits and diversification \\
Credit & PD/LGD/EAD, spread duration, covenant headroom & Issuer and correlated-exposure limits \\
Liquidity & Bid depth, days to liquidate, haircut & Liquidity ladder and stress cash buffer \\
Rates & Duration, DV01, convexity, curve shocks & Bucket/key-rate limits \\
FX/sovereign & Open currency and transfer exposure & Hedging, country limits, custody review \\
Concentration & Issuer/industry/driver/factor share & Multi-layer concentration limits \\
Counterparty & Exposure after collateral and wrong-way risk & Netting, collateral and thresholds \\
Model/governance & Validation exceptions, stale input, overrides & Independent validation and audit \\
\bottomrule\end{tabularx}\end{center}

\section{Portfolio exposure aggregation and look-through}
Measure economic exposure to issuer, ultimate parent, sector, geography, counterparties, liquidity buckets, and correlated cash-flow drivers. Equity and bonds of the same issuer are not independent positions. Treat look-through exposures in funds and vehicles when reliable; mark unknown exposures instead of assuming diversification.

\subsection{Concentration and sizing}
For position market value $MV_i$ and portfolio NAV,
\[
w_i=\frac{MV_i}{NAV},\quad HHI=\sum_i w_i^2,\quad N_{eff}=1/HHI
\]
for long-only weights summing to 1. These metrics reveal concentration but not common underlying risk drivers. \textbf{Illustrative policy, not universal advice:} single-issuer exposure requiring risk committee review above 10\% NAV, stress loss contribution above 2\% NAV, and cash-at-risk that jeopardizes known withdrawals. Real limits must be mandate-specific.

\section{Market and downside metrics}
\subsection{Volatility, covariance and expected shortfall}
For weights $\mathbf{w}$ and return covariance $\Sigma$, $\sigma_p=\sqrt{\mathbf{w}^\top\Sigma\mathbf{w}}$. With portfolio \emph{loss} random variable $L$, expected shortfall at confidence $\alpha$ can be defined as
\[
ES_\alpha=\frac{1}{1-\alpha}\int_\alpha^1 VaR_u(L)\,du.
\]
The integral handles continuous and discrete distributions using a quantile-integral convention; do not substitute an undefined arithmetic average for a discrete tail. VaR and ES rely on horizon, valuation, return distribution, and data regime. Report validation breaches and tail-sample limitations.

\subsection{Drawdown}
For NAV path $V_t$, define running peak $M_t=\max_{s\leq t}V_s$ and $DD_t=1-V_t/M_t$. Maximum drawdown $\max_t DD_t$ is a historical path statistic, not a worst-possible loss. Include recovery time and withdrawal obligations when deciding tolerability.

\section{Credit loss and fixed-income risk budgeting}
\[
EL_i=PD_i\,LGD_i\,EAD_i;\qquad
\Delta P\approx-D_{mod}P\Delta y+\frac12 Conv\,P(\Delta y)^2.
\]
Credit expected loss is only one layer. Distinguish maturity versus spread changes, correlated default clusters, recovery rates in downturns, covenant slippage, refinancing walls and cross-default effects. Reconcile bond market values, accrued interest, collateral and legal maturity dates.

\section{Liquidity engineering}
Model gross redemption/withdrawal schedule, legally restricted cash, collateral calls, bid--ask spread, position size relative to market depth, and time to exit. Define liquidity coverage over horizon $h$:
\[
LCR_h=\frac{\text{unencumbered stress-realizable assets by }h}{\text{net expected cash outflows by }h}.
\]
A ratio above 1 is necessary only within this simplified policy test; it does not establish adequacy if outflows, settlement or correlated market sales are mis-modeled. Apply tiered haircuts and exclude assets that cannot settle before payments fall due.

\section{Required stress scenarios}
\begin{enumerate}
\item \textbf{Recession + credit contraction:} earnings $-30\%$, financing rates +300bps, lower refinancing availability, recovery haircuts.
\item \textbf{Inflation/rate repricing:} nonparallel curve shift, pricing margin pressure, currency move, duration/convexity effects.
\item \textbf{Concentrated issuer failure:} simultaneous equity wipeout, debt restructuring, cross-guarantee litigation, exit liquidity freeze.
\item \textbf{Liquidity run:} large investor withdrawals, secondary-market discounts, derivatives margin calls, trading halts.
\item \textbf{International shock:} devaluation, capital controls, local-currency debt mismatch and custody/settlement obstruction.
\item \textbf{Model deception:} accounting restatement, previously unseen liability or governance fraud invalidating earning power.
\end{enumerate}

\section{Risk appetite to control mapping}
\begin{center}\small\begin{tabularx}{\textwidth}{p{0.26\textwidth}X X}\toprule Condition & Interpretation & Mandated action \\
\midrule Missing current exposure data & Risk cannot be bounded & Hold/stop incremental exposure \\
Hard mandate breach & Unauthorized position or class & Escalate, remediate by deadline \\
Funding shortfall scenario & Liquidity risk threatens solvency & Reduce commitments; increase reliable liquidity \\
Issuer thesis invalidated & Evidence contradicts underwriting & Re-underwrite; consider exit, do not auto-average \\
Significant model drift & Predictions and cash differ materially & Review model, challenge inputs, recalibrate \\
Concentrated adverse exposure & A single factor dominates tail & Hedge, reduce or authorize recorded exception \\
\bottomrule\end{tabularx}\end{center}

\section{Portfolio monitoring operating model}
\textbf{Daily:} prices, cash/settlement, exposures, limit alerts, liquidity, ratings and material corporate events. \textbf{Weekly:} exception queue, scenario changes, issuer concentration and upcoming maturities. \textbf{Monthly:} model validation, rolling drawdowns, cash budgets, hedge effectiveness and peer risk. \textbf{Quarterly:} issuer thesis restatement, investment-policy review, counterparty review, and benchmark-to-mandate reconciliation.

\section{Risk register template}
\small\begin{longtable}{p{0.12\textwidth}p{0.15\textwidth}p{0.19\textwidth}p{0.18\textwidth}p{0.20\textwidth}}
\toprule ID & Threat & Metric/trigger & Owner & Escalation/closure proof \\
\midrule\endhead
R-01 & Forced sale & Cash deficit at $h$ & Treasury/risk & Liquidity plan; verified cash \\
R-02 & Correlated failure & Concentrated stress loss & Portfolio manager & Reduced position; limit recheck \\
R-03 & Covenant deterioration & Headroom below policy floor & Credit analyst & Re-underwritten credit note \\
R-04 & Valuation uncertainty & Downside interval falls below price & Research lead & Updated assumptions and decision \\
R-05 & Stale or broken model & Source/reconciliation failure & Model owner & Reproducible validation pass \\
\bottomrule\end{longtable}\normalsize

\section{Validation and governance standards}
Separate model developer, independent challenger and approver. Freeze production model versions for completed decisions. Maintain data contracts, outlier checks, reproducible scenario seeds, retrospective forecast errors and explicit model risk. Do not use one confidence level or normal-distribution assumption as the sole indicator of a portfolio's failure modes. Escalate uncertainties that cannot be reliably quantified.

\end{document}
